\documentclass[11pt,a4paper]{article}

\usepackage[utf8]{inputenc}
\usepackage[T1]{fontenc}
\usepackage{amsmath,amssymb,amsthm}
\usepackage{bm}
\usepackage{geometry}
\usepackage{booktabs}
\usepackage{hyperref}
\usepackage{cleveref}

\geometry{margin=2.6cm}

\newtheorem{remark}{Remark}
\newtheorem{assumption}{Assumption}

\newcommand{\R}{\mathbb{R}}
\newcommand{\Xset}{\mathcal{X}}
\newcommand{\Uset}{\mathcal{U}}
\newcommand{\Qset}{\mathcal{Q}}
\newcommand{\Sset}{\mathcal{S}}
\newcommand{\vecx}{\bm{x}}
\newcommand{\vecu}{\bm{u}}
\newcommand{\vecw}{\bm{w}}
\newcommand{\vecc}{\bm{c}}
\newcommand{\vectheta}{\bm{\theta}}
\newcommand{\veceta}{\bm{\eta}}
\newcommand{\veclambda}{\bm{\lambda}}
\newcommand{\vecf}{\bm{f}}
\newcommand{\vecg}{\bm{g}}
\newcommand{\matG}{\bm{G}}

\title{Two Dimensional Vehicle Model}
\author{}
\date{}

% The capability vector has thirteen entries and the extended constraint
% vector eleven; amsmath's default matrix limit is ten, which silently
% wraps them into extra rows.
\setcounter{MaxMatrixCols}{24}

\begin{document}

\maketitle

\section{Overview}

The vehicle model is described by three complementary parts:

\begin{enumerate}
  \item a \emph{dynamics model}, which describes how the vehicle state evolves;
  \item a \emph{constraints model}, which defines admissible states and motion
        commands; and
  \item a \emph{capability model}, which describes what the vehicle can
        currently achieve.
\end{enumerate}

The separation is deliberate. The dynamics model is what the simulation core
integrates. The constraints model is a \emph{declaration} of the region in which
the dynamics model is claimed to be valid. The capability model is the
component's external interface: it answers questions about achievable motion
without requiring the dynamics to be integrated forward, and it is what guidance,
motion planning, and mission objective management consume.

The model is first introduced without discrete operating modes. A second
formulation then extends the model with two propulsion modes, nominal and boost.

\section{Conventions}
\label{sec:conventions}

\paragraph{Frame.}
A planar projection of the standard North\,--\,East\,--\,Down frame is used.
The coordinate $p_x$ is measured towards north and $p_y$ towards east, both in
metres. The vertical axis is suppressed; the vehicle is assumed to operate at
constant altitude.

\paragraph{Heading.}
The heading $\psi$ is measured clockwise from north, that is, from the positive
$p_x$ axis towards the positive $p_y$ axis. This is the conventional aviation
sense and is retained in preference to the mathematical convention.

\paragraph{Turn rate.}
The turn rate $\omega = \dot\psi$ is positive for a right turn, corresponding to
a positive bank angle $\varphi$.

\paragraph{Airspeed.}
The state $v$ is \emph{airspeed}. Wind enters the position kinematics only.
Consequently $\psi$ is an air relative heading, and the ground track differs from
$\psi$ whenever the wind is non zero. This distinction is retained because the
aerodynamic model is a function of airspeed while any external observer of the
vehicle measures ground referenced motion.

\paragraph{Units.}
SI base units throughout. Angles are in radians.

\section{Baseline Vehicle Model}

\subsection{Dynamics model}

The baseline continuous time dynamics are written as
\begin{equation}
  \dot{\vecx} = \vecf(\vecx, \vecu, \vectheta, \veceta) + \matG(\vecx)\vecw .
  \label{eq:baseline-dynamics}
\end{equation}

The state is
\begin{equation}
  \vecx = \begin{bmatrix} p_x & p_y & \psi & v & m \end{bmatrix}^{T} \in \R^{5},
\end{equation}
where $p_x$ and $p_y$ are position coordinates, $\psi$ is heading, $v$ is
airspeed, and $m$ is total vehicle mass.

The continuous motion command is
\begin{equation}
  \vecu = \begin{bmatrix} T & \omega \end{bmatrix}^{T} \in \R^{2},
\end{equation}
where $T$ is thrust and $\omega$ is turn rate.

The vehicle parameter vector is
\begin{equation}
  \vectheta = \begin{bmatrix} c_p & c_i & c_\ell & c_{\mathrm{TSFC}} \end{bmatrix}^{T},
  \label{eq:theta}
\end{equation}
where $c_p$ is the parasite drag parameter, $c_i$ is the induced drag parameter,
$c_\ell$ is the maximum lift parameter, and $c_{\mathrm{TSFC}}$ is the thrust
specific fuel consumption parameter.

The environmental parameter vector is
\begin{equation}
  \veceta = \begin{bmatrix} g & \rho \end{bmatrix}^{T},
\end{equation}
where $g$ is gravitational acceleration and $\rho$ is air density.

\subsubsection{Aerodynamic model}

The drag model follows the conventional parabolic drag polar
\cite{nasa-drag,nasa-induced},
\begin{equation}
  C_D = C_{D_0} + k C_L^2, \qquad k = \frac{1}{\pi e A\!R},
\end{equation}
where $C_D$ is total drag coefficient, $C_{D_0}$ is zero lift drag coefficient,
$C_L$ is lift coefficient, $A\!R$ is wing aspect ratio, and $e$ is the Oswald
efficiency factor.

Dynamic pressure is $\tfrac{1}{2}\rho v^2$, and lift and drag are written as
\begin{equation}
  L = \tfrac{1}{2}\rho v^{2} S C_L, \qquad
  D = \tfrac{1}{2}\rho v^{2} S C_D,
  \label{eq:lift-drag}
\end{equation}
where $S$ is the wing reference area \cite{nasa-lift}.

\paragraph{Parasite drag.}
The parasite drag is
\begin{equation}
  D_p = \tfrac{1}{2}\rho S C_{D_0} v^{2} .
\end{equation}
Defining
\begin{equation}
  c_p := \tfrac{1}{2} S C_{D_0},
  \label{eq:cp}
\end{equation}
gives
\begin{equation}
  D_p = \rho c_p v^{2} .
\end{equation}

\paragraph{Induced drag.}
For induced drag, substitution of $C_L$ from \cref{eq:lift-drag} gives
\begin{equation}
  D_i = \frac{2kL^{2}}{\rho S v^{2}} .
\end{equation}
In straight, steady, level flight $L = mg$, giving
\begin{equation}
  D_{i,0} = \frac{2k(mg)^{2}}{\rho S v^{2}} .
\end{equation}
Using $k = 1/(\pi e A\!R)$ and defining
\begin{equation}
  c_i := \frac{2g^{2}}{\pi e A\!R\, S},
  \label{eq:ci}
\end{equation}
the baseline induced drag becomes
\begin{equation}
  D_{i,0} = \frac{c_i}{\rho}\,\frac{m^{2}}{v^{2}} .
\end{equation}

\paragraph{Maximum lift.}
The largest lift the vehicle can generate at a given airspeed follows from
\cref{eq:lift-drag} evaluated at the maximum lift coefficient $C_{L,\max}$.
Defining
\begin{equation}
  c_\ell := \tfrac{1}{2} S C_{L,\max},
  \label{eq:cl}
\end{equation}
in direct analogy with \cref{eq:cp}, gives
\begin{equation}
  L_{\max}(v) = \rho c_\ell v^{2} .
  \label{eq:lmax}
\end{equation}

\begin{remark}
The three lumped parameters $c_p$, $c_i$ and $c_\ell$ absorb the geometry
$S$, $C_{D_0}$, $e$, $A\!R$ and $C_{L,\max}$, so that a vehicle may be specified
either by the lumped parameters directly or by the underlying geometry.
\end{remark}

\subsubsection{Coordinated turn}

During a coordinated level turn with bank angle $\varphi$, the vertical
component of lift balances weight,
\begin{equation}
  L\cos\varphi = mg,
\end{equation}
so the load factor is
\begin{equation}
  n = \frac{L}{mg} = \frac{1}{\cos\varphi} .
  \label{eq:load-factor-def}
\end{equation}
The horizontal lift component provides the centripetal force,
\begin{equation}
  L\sin\varphi = \frac{mv^{2}}{R},
\end{equation}
where $R$ is turn radius. Since $\omega = v/R$,
\begin{equation}
  \omega = \frac{g\tan\varphi}{v},
  \label{eq:omega-phi}
\end{equation}
and consequently
\begin{equation}
  n^{2} = 1 + \left(\frac{v\omega}{g}\right)^{2} .
  \label{eq:n-omega}
\end{equation}
The sign convention of \cref{sec:conventions} is consistent with
\cref{eq:omega-phi}: a positive bank angle produces a positive turn rate and
hence a right turn.

\subsubsection{Total drag}

Because induced drag is proportional to $L^{2}$, it scales with $n^{2}$. Using
\cref{eq:n-omega}, the total drag is
\begin{equation}
  D(v, m, \omega; \vectheta, \veceta)
  = \rho c_p v^{2}
  + \frac{c_i}{\rho}\,\frac{m^{2}}{v^{2}}
    \left[ 1 + \left(\frac{v\omega}{g}\right)^{2} \right] .
  \label{eq:drag}
\end{equation}
It is convenient to name the induced drag coefficient at unit load factor,
\begin{equation}
  A(v,m) := \frac{c_i m^{2}}{\rho v^{2}},
  \label{eq:A}
\end{equation}
so that \cref{eq:drag} may be written compactly as
\begin{equation}
  D = \rho c_p v^{2} + A(v,m)\left[1 + \left(\frac{v\omega}{g}\right)^{2}\right].
  \label{eq:drag-compact}
\end{equation}

\subsubsection{Nominal dynamics}

The nominal dynamics function is
\begin{equation}
  \vecf(\vecx, \vecu, \vectheta, \veceta) =
  \begin{bmatrix}
    v\cos\psi \\[2pt]
    v\sin\psi \\[2pt]
    \omega \\[2pt]
    \dfrac{T - D(v, m, \omega; \vectheta, \veceta)}{m} \\[8pt]
    -c_{\mathrm{TSFC}} T
  \end{bmatrix}.
  \label{eq:f-nominal}
\end{equation}

\begin{remark}
Under the convention of \cref{sec:conventions}, with $p_x$ towards north and
$\psi$ clockwise from north, the position kinematics retain the same algebraic
form as under the mathematical convention. Only the interpretation of the axes
and the sign of $\omega$ differ.
\end{remark}

A compact disturbance vector is
\begin{equation}
  \vecw = \begin{bmatrix} w_x & w_y & w_\psi & F_v \end{bmatrix}^{T},
\end{equation}
where $w_x$ and $w_y$ are wind velocity components, $w_\psi$ is a turn rate
disturbance, and $F_v$ is a signed longitudinal disturbance force.

The disturbance matrix is
\begin{equation}
  \matG(\vecx) =
  \begin{bmatrix}
    1 & 0 & 0 & 0 \\
    0 & 1 & 0 & 0 \\
    0 & 0 & 1 & 0 \\
    0 & 0 & 0 & \tfrac{1}{m} \\
    0 & 0 & 0 & 0
  \end{bmatrix}.
\end{equation}

The complete baseline dynamics are therefore
\begin{align}
  \dot p_x &= v\cos\psi + w_x, \\
  \dot p_y &= v\sin\psi + w_y, \\
  \dot\psi &= \omega + w_\psi, \\
  \dot v   &= \frac{T - D(v, m, \omega; \vectheta, \veceta) + F_v}{m}, \\
  \dot m   &= -c_{\mathrm{TSFC}} T .
\end{align}

\subsection{Constraints model}

The baseline constraint parameters are collected in
\begin{equation}
  \veclambda = \begin{bmatrix}
    T_{\min} & T_{\max} & n_{\max} & \omega_{\mathrm{cap}} &
    v_{\min} & v_{\max} & m_{\mathrm{dry}} & m_{\max}
  \end{bmatrix}^{T},
  \label{eq:lambda}
\end{equation}
where $n_{\max}$ is the structural load factor limit and $\omega_{\mathrm{cap}}$
is an absolute turn rate bound representing roll and control authority.

\subsubsection{Turn rate limit}

The turn rate is not bounded by a constant. Two distinct limits act, and the
binding one depends on the operating point.

The \emph{lift limited} load factor follows from \cref{eq:lmax} and
\cref{eq:load-factor-def},
\begin{equation}
  n_\ell(v, m) = \frac{L_{\max}(v)}{mg} = \frac{\rho c_\ell v^{2}}{mg},
  \label{eq:n-lift}
\end{equation}
and the available load factor is the smaller of the lift limit and the
structural limit,
\begin{equation}
  n_{\mathrm{av}}(v, m) = \min\{\, n_{\max},\; n_\ell(v, m) \,\}.
  \label{eq:n-available}
\end{equation}
Inverting \cref{eq:n-omega} gives the corresponding turn rate bound,
\begin{equation}
  \omega_{\max}(v, m) =
  \begin{cases}
    \min\left\{ \omega_{\mathrm{cap}},\;
      \dfrac{g}{v}\sqrt{n_{\mathrm{av}}^{2}(v,m) - 1} \right\},
      & n_{\mathrm{av}}(v,m) > 1, \\[10pt]
    0, & \text{otherwise.}
  \end{cases}
  \label{eq:omega-max}
\end{equation}

\begin{remark}
The lift limit dominates at low airspeed and the structural limit at high
airspeed. The crossover, where $n_\ell(v,m) = n_{\max}$, is the corner speed
defined in \cref{eq:v-corner}, at which the instantaneous turn rate attains its
maximum. A constant turn rate bound would imply an unbounded load factor as
airspeed increases and would permit turns the wing cannot support at low
airspeed.
\end{remark}

\subsubsection{Admissible sets}

The admissible input set is
\begin{equation}
  \Uset(\vecx, \veclambda) = \left\{
    \vecu \in \R^{2} : T_{\min} \le T \le T_{\max},\;
    |\omega| \le \omega_{\max}(v, m)
  \right\}.
  \label{eq:U}
\end{equation}
The admissible state set is
\begin{equation}
  \Xset(\veclambda) = \left\{
    \vecx \in \R^{5} : v_{\mathrm{s}}(m,1) \le v \le v_{\max},\;
    v \ge v_{\min},\; m_{\mathrm{dry}} \le m \le m_{\max}
  \right\},
  \label{eq:X}
\end{equation}
where $v_{\mathrm{s}}$ is the stall speed defined in \cref{eq:v-stall}.

\begin{remark}[The mass ceiling is a loading constraint, not a flight one]
$m \le m_{\max}$ is the only element of $\Xset$ that cannot be violated by
how the vehicle is flown. Mass falls monotonically as fuel burns, so a
trajectory beginning inside the set stays inside it, and no control input can
push it out. It is stated here because it is a property of the airframe and
belongs with the airframe's other limits, and because it gives a scenario that
hand-builds an overloaded initial state the same finding that a badly
specified platform receives from the composition-time mass budget. Enforcing
it is a composition concern; declaring it is the vehicle's.
\end{remark}

\begin{remark}[Enforcement is delegated]
\label{rem:enforcement}
The constraints model \emph{declares} the region in which the dynamics model is
claimed to be valid. It does not specify what occurs when a command or a state
lies outside that region, and the vehicle model performs no projection,
saturation or clamping of its own.

Responsibility for respecting $\Uset$ and $\Xset$ lies with the guidance and
motion planning logic, possibly in combination with a runtime assurance
mechanism interposed between the planner and the vehicle. Placing enforcement in
the vehicle would conceal control design faults, would make the enforcement
policy an unstated property of the vehicle rather than an explicit design
decision, and would prevent different enforcement strategies from being compared
against the same vehicle.

Outside $\Uset$ and $\Xset$ the dynamics remain integrable but the model is not
claimed to be a faithful representation of the vehicle.
\end{remark}

\begin{remark}
The input set $\Uset$ is state dependent, since $\omega_{\max}$ is a function of
airspeed and mass. The vehicle exposes $\omega_{\max}(v,m)$ through the
capability model, so that guidance logic may evaluate the bound at the current
operating point without reproducing the aerodynamic model.
\end{remark}

\subsection{Capability model}

The capability model is written generally as
\begin{equation}
  \vecc = \vecg(\vecx, \vectheta, \veceta, \veclambda, \vecw).
  \label{eq:capability-general}
\end{equation}

A capability vector, evaluated at a specified turn rate $\omega$, is
\begin{equation}
  \vecc = \begin{bmatrix}
    T_{\mathrm{av}} & T_{\mathrm{req}} &
    a_{\max} & a_{\min} &
    n_{\mathrm{av}} & \omega_{\mathrm{av}} & \omega_{\mathrm{sus}} & R_{\min} &
    v_{\mathrm{s}} & v_{\mathrm{c}} &
    m_{\mathrm{fuel}} & t_{\mathrm{end}} & r_{\mathrm{sp}}
  \end{bmatrix}^{T}.
  \label{eq:c}
\end{equation}

\paragraph{Thrust.}
The available thrust is
\begin{equation}
  T_{\mathrm{av}} =
  \begin{cases}
    T_{\max}, & m > m_{\mathrm{dry}}, \\
    0, & m \le m_{\mathrm{dry}} .
  \end{cases}
\end{equation}
The thrust required to hold the current airspeed at the specified turn rate is
\begin{equation}
  T_{\mathrm{req}}(v, m, \omega) = D(v, m, \omega; \vectheta, \veceta).
  \label{eq:t-req}
\end{equation}

\paragraph{Longitudinal acceleration.}
For a specified turn rate $\omega$, the instantaneous maximum and minimum
longitudinal accelerations are
\begin{align}
  a_{\max} &= \frac{T_{\mathrm{av}} - D(v, m, \omega; \vectheta, \veceta) + F_v}{m}, \\
  a_{\min} &= \frac{T_{\min} - D(v, m, \omega; \vectheta, \veceta) + F_v}{m}.
\end{align}

\paragraph{Instantaneous turn performance.}
The available load factor $n_{\mathrm{av}}$ is given by \cref{eq:n-available} and
the available turn rate by
\begin{equation}
  \omega_{\mathrm{av}}(v, m) = \omega_{\max}(v, m),
\end{equation}
with the corresponding minimum turn radius
\begin{equation}
  R_{\min} = \frac{v}{\omega_{\mathrm{av}}}, \qquad \omega_{\mathrm{av}} > 0 .
\end{equation}

\paragraph{Sustained turn performance.}
The available turn rate is instantaneous: holding it in general requires drag to
exceed available thrust, so airspeed decays. The \emph{sustained} turn rate is
the largest turn rate at which thrust and drag balance, so that $\dot v = 0$ in
the absence of disturbance. Setting $T_{\mathrm{av}} = D(v, m, \omega)$ in
\cref{eq:drag-compact} and solving for $\omega$ gives the closed form
\begin{equation}
  \omega_{\mathrm{sus}}(v, m) =
  \min\left\{
    \omega_{\mathrm{av}}(v,m),\;
    \frac{g}{v}
    \sqrt{\frac{T_{\mathrm{av}} - \rho c_p v^{2} - A(v,m)}{A(v,m)}}
  \right\},
  \label{eq:omega-sustained}
\end{equation}
defined whenever $T_{\mathrm{av}} > \rho c_p v^{2} + A(v,m)$, and zero otherwise.
The condition for a non zero sustained turn rate is precisely the condition that
the vehicle can maintain wings level flight at the given airspeed.

\begin{remark}
The interval $[\omega_{\mathrm{sus}}, \omega_{\mathrm{av}}]$ is the region in
which the vehicle may turn only by expending kinetic energy. The width of this
interval is the quantity a motion planner trades against, and is therefore
exposed explicitly rather than left to be inferred.
\end{remark}

\paragraph{Characteristic speeds.}
The stall speed at load factor $n$ follows from inverting \cref{eq:n-lift},
\begin{equation}
  v_{\mathrm{s}}(m, n) = \sqrt{\frac{n m g}{\rho c_\ell}},
  \label{eq:v-stall}
\end{equation}
and the corner speed, at which the lift and structural limits coincide and the
instantaneous turn rate is maximised, is
\begin{equation}
  v_{\mathrm{c}}(m) = v_{\mathrm{s}}(m, 1)\sqrt{n_{\max}}
  = \sqrt{\frac{n_{\max} m g}{\rho c_\ell}} .
  \label{eq:v-corner}
\end{equation}

\paragraph{Fuel and endurance.}
Remaining fuel mass is
\begin{equation}
  m_{\mathrm{fuel}} = \max\{m - m_{\mathrm{dry}},\, 0\} .
\end{equation}
At the thrust required for steady flight, an instantaneous endurance estimate is
\begin{equation}
  t_{\mathrm{end}} = \frac{m_{\mathrm{fuel}}}{c_{\mathrm{TSFC}} T_{\mathrm{req}}},
  \qquad T_{\mathrm{req}} > 0,
  \label{eq:endurance}
\end{equation}
and the corresponding specific range, that is, distance flown per unit mass of
fuel, is
\begin{equation}
  r_{\mathrm{sp}} = \frac{v}{c_{\mathrm{TSFC}} T_{\mathrm{req}}} .
  \label{eq:specific-range}
\end{equation}

\begin{remark}
The speed for maximum endurance minimises $T_{\mathrm{req}}(v,m,0)$ and the
speed for maximum range maximises $v / T_{\mathrm{req}}(v,m,0)$. Both follow from
\cref{eq:drag} by differentiation and are available to mission objective
management as derived quantities.
\end{remark}

\begin{remark}[True and believed capability]
\label{rem:true-believed}
Evaluating \cref{eq:capability-general} requires the disturbance $\vecw$, which
is a property of the environment and is not observable by the vehicle. Two
distinct evaluations are therefore meaningful:
\begin{equation}
  \vecc = \vecg(\vecx, \vectheta, \veceta, \veclambda, \vecw),
  \qquad
  \hat{\vecc} = \vecg(\hat{\vecx}, \vectheta, \veceta, \veclambda, \hat{\vecw}),
\end{equation}
where $\hat{\vecx}$ and $\hat{\vecw}$ denote estimates. The former is the true
capability and is available only to the simulation. The latter is the capability
the vehicle believes it has, and is what guidance and planning must use. The
discrepancy between the two is a property of the estimation problem and is
intentionally preserved rather than removed.
\end{remark}

\section{Vehicle Model with Nominal and Boost Modes}

\subsection{Mode dependent dynamics}

The vehicle is extended with the discrete propulsion mode
\begin{equation}
  q \in \Qset, \qquad \Qset = \{\mathrm{nom},\, \mathrm{boost}\} .
\end{equation}

The continuous motion command is unchanged. The continuous state is augmented
with a scalar \emph{thermal state} $s$, introduced in
\cref{sec:mode-switching}, which accumulates while boost is engaged:
\begin{equation}
  \vecx = \begin{bmatrix} p_x & p_y & \psi & v & m & s \end{bmatrix}^{T},
  \qquad
  \vecu = \begin{bmatrix} T & \omega \end{bmatrix}^{T} .
\end{equation}

The switched dynamics are written as
\begin{equation}
  \dot{\vecx} = \vecf_q(\vecx, \vecu, \vectheta, \veceta) + \matG(\vecx)\vecw,
  \qquad q \in \Qset .
\end{equation}

The parameter vector becomes
\begin{equation}
  \vectheta = \begin{bmatrix}
    c_p & c_i & c_\ell &
    c^{\mathrm{nom}}_{\mathrm{TSFC}} & c^{\mathrm{boost}}_{\mathrm{TSFC}} &
    \tau_{\mathrm{h}} & \tau_{\mathrm{c}}
  \end{bmatrix}^{T},
\end{equation}
with
\begin{equation}
  c^{\mathrm{boost}}_{\mathrm{TSFC}} > c^{\mathrm{nom}}_{\mathrm{TSFC}},
\end{equation}
where $\tau_{\mathrm{h}}$ and $\tau_{\mathrm{c}}$ are the thermal accumulation
and recovery time constants.

The mode dependent fuel consumption parameter is
\begin{equation}
  c_{\mathrm{TSFC},q} =
  \begin{cases}
    c^{\mathrm{nom}}_{\mathrm{TSFC}}, & q = \mathrm{nom}, \\
    c^{\mathrm{boost}}_{\mathrm{TSFC}}, & q = \mathrm{boost},
  \end{cases}
\end{equation}
and the mode dependent dynamics are
\begin{equation}
  \vecf_q =
  \begin{bmatrix}
    v\cos\psi \\[2pt]
    v\sin\psi \\[2pt]
    \omega \\[2pt]
    \dfrac{T - D(v, m, \omega; \vectheta, \veceta)}{m} \\[8pt]
    -c_{\mathrm{TSFC},q} T \\[4pt]
    \sigma_q(s)
  \end{bmatrix},
  \qquad
  \sigma_q(s) =
  \begin{cases}
    \dfrac{1}{\tau_{\mathrm{h}}}, & q = \mathrm{boost}, \\[8pt]
    -\dfrac{s}{\tau_{\mathrm{c}}}, & q = \mathrm{nom} .
  \end{cases}
  \label{eq:f-mode}
\end{equation}

The drag model and the disturbance model are unchanged. Boost mode increases
fuel consumption per unit thrust and, through the constraints model, permits a
larger thrust command.

\subsection{Mode dependent constraints model}

The constraint parameters are extended with mode dependent propulsion limits:
\begin{equation}
  \veclambda = \begin{bmatrix}
    T_{\min} & T^{\mathrm{nom}}_{\max} & T^{\mathrm{boost}}_{\max} &
    n_{\max} & \omega_{\mathrm{cap}} &
    v_{\min} & v^{\mathrm{nom}}_{\max} & v^{\mathrm{boost}}_{\max} &
    m_{\mathrm{dry}} & m_{\max} & m_{\mathrm{res}} & s_{\max}
  \end{bmatrix}^{T},
\end{equation}
with
\begin{equation}
  T^{\mathrm{boost}}_{\max} > T^{\mathrm{nom}}_{\max},
  \qquad
  v^{\mathrm{boost}}_{\max} > v^{\mathrm{nom}}_{\max},
\end{equation}
where $m_{\mathrm{res}}$ is a fuel reserve below which boost is inhibited and
$s_{\max}$ is the thermal limit.

The mode dependent input set is
\begin{equation}
  \Uset_q(\vecx, \veclambda) = \left\{
    \vecu : T_{\min} \le T \le T^{q}_{\max},\;
    |\omega| \le \omega_{\max}(v, m)
  \right\},
\end{equation}
and the mode dependent state set is
\begin{equation}
  \Xset_q(\veclambda) = \left\{
    \vecx : v_{\mathrm{s}}(m,1) \le v \le v^{q}_{\max},\;
    v \ge v_{\min},\;
    m_{\mathrm{dry}} \le m \le m_{\max},\;
    0 \le s \le s_{\max}
  \right\}.
\end{equation}

\begin{remark}[The turn rate bound is mode independent]
\label{rem:mode-independent-turn}
The bound $\omega_{\max}(v,m)$ of \cref{eq:omega-max} does not carry a mode
index. The instantaneous turn rate is limited by available lift and by structural
strength, neither of which is affected by propulsion setting. An earlier
formulation in which $\omega^{\mathrm{boost}}_{\max} > \omega^{\mathrm{nom}}_{\max}$
is not consistent with the aerodynamic model.

Boost nonetheless improves turn performance, but it does so through the
\emph{sustained} turn rate \cref{eq:omega-sustained}, which depends on available
thrust and is therefore mode dependent. This is the physically correct channel:
boost does not allow a tighter turn, it allows a given turn to be held without
losing airspeed.
\end{remark}

\subsection{Mode switching constraints}
\label{sec:mode-switching}

The mode is a discrete control input, but transitions between modes are not
assumed to be unrestricted. Let $q$ denote the current mode and $q^{+}$ the mode
after a switching decision. Admissible transitions are described by
\begin{equation}
  q^{+} \in \Sset_q(\vecx, \veclambda) \subseteq \Qset .
  \label{eq:switching-abstract}
\end{equation}

A concrete instantiation sufficient for the present model is
\begin{equation}
  \Sset_{q}(\vecx, \veclambda) =
  \begin{cases}
    \{\mathrm{nom}\},
      & s \ge s_{\max}
        \;\text{ or }\;
        m \le m_{\mathrm{dry}} + m_{\mathrm{res}}, \\[4pt]
    \{q\},
      & t - t_{q} < \Delta t_{\mathrm{dwell}}, \\[4pt]
    \{\mathrm{nom}, \mathrm{boost}\},
      & \text{otherwise,}
  \end{cases}
  \label{eq:switching-concrete}
\end{equation}
where $t_q$ is the time of the most recent mode transition and
$\Delta t_{\mathrm{dwell}}$ is a minimum dwell time. The cases are evaluated
in the order given.

\begin{remark}[Being forced out and being locked in are different]
An earlier formulation of \cref{eq:switching-concrete} listed the two
permissive cases by current mode and returned $\{\mathrm{nom}\}$ otherwise.
That conflates two situations which happen to share a branch. Exhausting the
thermal or fuel budget must force the aircraft \emph{out} of boost; not having
served the dwell must \emph{keep} it where it is. Under the earlier
formulation, $q = \mathrm{boost}$ with the dwell unserved returned
$\{\mathrm{nom}\}$, so boost was granted on one step and revoked on the next
indefinitely --- an anti-chattering condition that produced chattering. It was
found by simulating the model rather than by reading it, and the thermal state
never rose above two per cent of its limit.

The ordering is also load-bearing. Were the dwell tested first, an aircraft
reaching $s_{\max}$ shortly after engaging would be held in boost past its own
thermal limit, and the guard against chattering would defeat the guard against
destroying the engine.
\end{remark}

Three distinct effects are captured. The thermal state $s$, evolving according to
\cref{eq:f-mode}, limits the duration of continuous boost and imposes a recovery
period afterwards. The fuel reserve $m_{\mathrm{res}}$ prevents boost from being
selected when the remaining fuel is committed to recovery. The dwell time
$\Delta t_{\mathrm{dwell}}$ prevents chattering, which is otherwise a hazard both
for the planner and for numerical integration of the switched system.

Only the first of the three is physical. An engine genuinely cannot sustain
boost past its thermal limit, whereas an aircraft can light its afterburner on
the last hundred kilograms of fuel and can switch modes twice in a second; it
simply should not. The fuel reserve and the dwell time are therefore declared
policy, carried in $\veclambda$ alongside $v_{\min}$, which is likewise a
declared floor rather than an aerodynamic one. The distinction matters to
anyone configuring a different aircraft: the thermal limit is the entry they
cannot simply relax.

\begin{remark}[The switching set is declared, not enforced]
The mode is a discrete input, so \cref{eq:switching-abstract} is an input
constraint, and the declaration/enforcement separation of the constraints model
applies to it unchanged: the vehicle \emph{declares} $\Sset_q$ and does not
enforce it. A
projection is offered, and because a mode cannot be clipped by degree the way a
thrust command can, that projection is a fallback to $\mathrm{nom}$ rather than
a partial reduction. A commanded mode outside $\Sset_q$ that is flown anyway
remains integrable, is not claimed valid, and shows itself through
$\Xset_q$: $s$ runs past $s_{\max}$ and is reported. This is the same treatment
a thrust command above $T^{q}_{\max}$ receives.
\end{remark}

\begin{remark}[The dwell condition is an argument, not a state]
$t - t_{q}$ is history rather than state, so a declaration written as a pure
function of $\vecx$ and $\veclambda$ cannot evaluate it. Two resolutions are
available: carry $t_{q}$ as an additional state, or supply the elapsed time to
the declaration as an argument. The latter is preferred here. Adding a state to
the vector would place a decision-hygiene quantity inside the integrator and
the Jacobian, where it participates in nothing, and the party that selects the
mode necessarily observes its own transitions and can supply the interval at no
cost.
\end{remark}

\begin{remark}
\Cref{eq:switching-concrete} is one instantiation of the abstract relation
\cref{eq:switching-abstract} and is not intended to be exhaustive. Further
restrictions related to engine condition, operating point, or externally imposed
rules of engagement may be introduced without altering the structure of the
model, provided they are expressible as a function of $\vecx$ and $\veclambda$.
\end{remark}

\subsection{Mode dependent capability model}

The capability model becomes
\begin{equation}
  \vecc = \vecg_q(\vecx, \vectheta, \veceta, \veclambda, \vecw),
\end{equation}
with the following quantities acquiring a mode index.

The currently available thrust is
\begin{equation}
  T_{\mathrm{av},q} =
  \begin{cases}
    T^{\mathrm{nom}}_{\max}, & q = \mathrm{nom},\; m > m_{\mathrm{dry}}, \\
    T^{\mathrm{boost}}_{\max}, & q = \mathrm{boost},\; m > m_{\mathrm{dry}}, \\
    0, & m \le m_{\mathrm{dry}} .
  \end{cases}
\end{equation}
The maximum longitudinal acceleration follows as
\begin{equation}
  a_{\max,q} = \frac{T_{\mathrm{av},q} - D(v, m, \omega; \vectheta, \veceta) + F_v}{m},
\end{equation}
and the sustained turn rate from \cref{eq:omega-sustained} evaluated with
$T_{\mathrm{av},q}$,
\begin{equation}
  \omega_{\mathrm{sus},q}(v, m) =
  \min\left\{
    \omega_{\mathrm{av}}(v,m),\;
    \frac{g}{v}
    \sqrt{\frac{T_{\mathrm{av},q} - \rho c_p v^{2} - A(v,m)}{A(v,m)}}
  \right\}.
\end{equation}
By \cref{rem:mode-independent-turn}, the instantaneous quantities
$n_{\mathrm{av}}$, $\omega_{\mathrm{av}}$, $R_{\min}$, $v_{\mathrm{s}}$ and
$v_{\mathrm{c}}$ are mode independent.

The endurance estimate becomes
\begin{equation}
  t_{\mathrm{end},q} =
  \frac{m_{\mathrm{fuel}}}{c_{\mathrm{TSFC},q} T_{\mathrm{req}}},
  \qquad T_{\mathrm{req}} > 0 .
\end{equation}

The capability model also exposes whether boost is currently selectable, and for
how long it could be held,
\begin{equation}
  c_{\mathrm{boost}} =
  \begin{cases}
    1, & \mathrm{boost} \in \Sset_q(\vecx, \veclambda), \\
    0, & \mathrm{boost} \notin \Sset_q(\vecx, \veclambda),
  \end{cases}
  \qquad
  t_{\mathrm{boost}} = \tau_{\mathrm{h}}\,(s_{\max} - s),
  \label{eq:boost-available}
\end{equation}
where $t_{\mathrm{boost}}$ is the time remaining before the thermal limit is
reached at continuous boost. Both quantities are required by mission objective
management, which must reason about whether a boost committed now remains
available later.

\section{Compact Summary}

The baseline model is
\begin{equation}
  \begin{aligned}
    \dot{\vecx} &= \vecf(\vecx, \vecu, \vectheta, \veceta) + \matG(\vecx)\vecw, \\
    \vecx &\in \Xset(\veclambda), \\
    \vecu &\in \Uset(\vecx, \veclambda), \\
    \vecc &= \vecg(\vecx, \vectheta, \veceta, \veclambda, \vecw).
  \end{aligned}
\end{equation}

The mode dependent extension is
\begin{equation}
  \begin{aligned}
    \dot{\vecx} &= \vecf_q(\vecx, \vecu, \vectheta, \veceta) + \matG(\vecx)\vecw,
      \qquad q \in \Qset, \\
    \vecx &\in \Xset_q(\veclambda), \\
    \vecu &\in \Uset_q(\vecx, \veclambda), \\
    q^{+} &\in \Sset_q(\vecx, \veclambda), \\
    \vecc &= \vecg_q(\vecx, \vectheta, \veceta, \veclambda, \vecw).
  \end{aligned}
\end{equation}

In both cases the membership relations are declarations of validity rather than
enforced projections, in the sense of \cref{rem:enforcement}.

\section{Assumptions and Limitations}

The model assumes a planar point mass vehicle, a coordinated steady level turn,
negligible sideslip, negligible altitude change, negligible roll transients,
instantaneous adjustment of lift, and validity of the parabolic drag polar within
the considered operating region. The coordinated turn relations are standard
results in aircraft flight mechanics \cite{faa-phak}.

Several consequences follow from the planar assumption and are recorded
explicitly:

\begin{enumerate}
  \item Potential energy is absent. Manoeuvres that trade altitude for airspeed,
        or the converse, cannot be represented, and the energy state of the
        vehicle is fully described by airspeed.
  \item Air density $\rho$ is a fixed environmental parameter rather than a
        function of altitude, so the model describes a single flight level.
  \item Thrust specific fuel consumption is constant within a mode, whereas in
        practice it depends on airspeed and altitude.
  \item Turn rate and thrust are commanded directly, so roll response and engine
        spool dynamics are absent. A first order lag on either command,
        $\dot\omega = (\omega_{\mathrm{c}} - \omega)/\tau_\omega$ and
        $\dot T = (T_{\mathrm{c}} - T)/\tau_T$, may be introduced by augmenting
        the state without altering the remainder of the model.
  \item The induced drag model degrades as airspeed approaches the stall speed
        \cref{eq:v-stall}, and the model is not claimed to be valid below it.
\end{enumerate}

\begin{thebibliography}{9}

\bibitem{nasa-drag}
NASA Glenn Research Center,
\emph{Drag Coefficient}, Beginner's Guide to Aeronautics.
\url{https://www1.grc.nasa.gov/beginners-guide-to-aeronautics/drag-coefficient/}

\bibitem{nasa-induced}
NASA Glenn Research Center,
\emph{Induced Drag Coefficient}, Beginner's Guide to Aeronautics.
\url{https://www1.grc.nasa.gov/beginners-guide-to-aeronautics/induced-drag-coefficient/}

\bibitem{nasa-lift}
NASA Glenn Research Center,
\emph{Lift Equation}, Beginner's Guide to Aeronautics.
\url{https://www1.grc.nasa.gov/beginners-guide-to-aeronautics/lift-equation/}

\bibitem{faa-phak}
Federal Aviation Administration,
\emph{Pilot's Handbook of Aeronautical Knowledge}, FAA-H-8083-25C,
Chapter 5: Aerodynamics of Flight, 2023.
\url{https://www.faa.gov/sites/faa.gov/files/07_phak_ch5_0.pdf}

\end{thebibliography}

\end{document}
